\section{Canonical factor order and an enrichment hierarchy}
\label{sec:enrichment}

\subsection{Factorial search is unnecessary}

The annihilator of an exponential polynomial supplies a multiset of factors,
but initially seems to leave their order to search. The initial coordinates
change with the order, so arbitrary ordering can indeed fail. Fortunately,
there is a best order in the relevant sense.

\begin{theorem}[Ascending-order dominance]
\label{thm:order}
Fix a function, an anchor, and a multiset of constant real differential factors
whose product annihilates the function. If any ordering of these factors gives
nonnegative ladder initial values, then their nondecreasing ordering does too.
\end{theorem}
\begin{proof}
Consider an adjacent inversion $\alpha>\beta$ after a common prefix has produced
$q$. In the original order the intermediate function is
$u=(\D-\alpha)q$. In the swapped order it is
\[
 \widetilde u=(\D-\beta)q=u+(\alpha-\beta)q.
\]
The original successful ladder gives $q(a)\ge0$ and $u(a)\ge0$, hence
$\widetilde u(a)\ge0$. After both factors the functions are exactly equal,
since $(\D-\alpha)$ and $(\D-\beta)$ commute. Every earlier and later initial
value is unchanged, as is the terminal zero. Each adjacent inversion can
therefore be removed without losing acceptance. Bubble sorting finishes.
\end{proof}

\begin{corollary}
For the native multiset containing each rate $\lambda$ exactly $d_\lambda+1$
times, checking one ascending ladder decides whether any permutation of that
multiset admits a ladder certificate, provided the initial signs are determined
exactly.
\end{corollary}

The qualification is important. This is completeness for a \emph{certificate
class}, not for positivity of exponential polynomials. At the origin, the
prototype's rational signs are exact, so it implements precisely that
fixed-multiset test. At other anchors, its fixed-precision enclosures can fail
to establish a true nonnegative initial value. The ordering theorem remains
true, but bounded numerical sign discovery is an additional source of
\unknown{}.

With $N$ native factors and at most $S\le N$ coefficients in the corresponding
dense block representation, a straightforward implementation uses $O(NS)$
rational arithmetic operations. Sorting rates costs $O(|\Lambda|\log|\Lambda|)$
comparisons; repeated rates need not actually be sorted one copy at a time.
Numerator and denominator growth must be accounted for separately. This is
not a bit-complexity or timing claim.

The implementation should therefore never offer a large permutation budget for
this lane. A failed ascending test is a reason to enrich the language, move the
anchor, split the domain, factor the target, or route elsewhere. Reordering the
same factors cannot reveal a successful certificate that ascending order missed
at exact signs.

\subsection{Adding factors without changing the source function}

A fixed multiset can be too weak. Consider
\begin{equation}
 f_\varepsilon(x)=e^{2x}-4e^x+4+\varepsilon,
 \qquad \varepsilon>0.
 \label{eq:epsilon-family}
\end{equation}
The native factors are $(0,1,2)$. The first step tests $f_\varepsilon(0)=1+\varepsilon>0$,
but its successor has value $f_\varepsilon'(0)=-2$. Thus the native ladder
fails, regardless of factor order by Theorem~\ref{thm:order}. Nevertheless
$f_\varepsilon=(e^x-2)^2+\varepsilon$ is strictly positive.

If $A(\D)f=0$, then $(\D-\beta)A(\D)f=0$ for every real $\beta$.
We can therefore add factors without changing the input function. They change
the coordinates used to establish positivity, not the theorem to be proved.

\begin{proposition}[Monotonicity under multiset extension]
\label{prop:extension}
If a function has an accepted ascending ladder for a multiset $M$, then it has
an accepted ascending ladder for every finite multiset containing $M$.
\end{proposition}
\begin{proof}
Append the additional factors after the existing terminal zero. Every new
intermediate function and initial value is zero, so the extended unsorted
ladder is accepted. Apply Theorem~\ref{thm:order} to sort the full multiset.
\end{proof}

The enlarged cones can be strictly larger. For $f_1=e^{2x}-4e^x+5$, four
additional factors equal to $-4$ give
\[
 (-4,-4,-4,-4,0,1,2),
\]
with initial vector
\[
 (2,6,16,36,76,92,2592).
\]
Every entry is positive and the final function is zero. This is a completely
rational certificate of strict positivity on $[0,\infty)$.

The prototype calls these \emph{ghost factors}; the article generally calls
them auxiliary factors. They are not additional program variables or new
continuous state variables. In particular, the term is not a claim that this
construction implements the differential-ghost rule of a general dynamic
logic~\cite{differential-invariance}.

\subsection{A bounded but useful search strategy}

For a native minimum rate $\lambda_{\min}$, try
\[
 \beta\in\{\lambda_{\min}-1,\lambda_{\min}-2,
           \lambda_{\min}-4,\ldots,\lambda_{\min}-64\}.
\]
For each $\beta$, append repeated copies \emph{before} the native ascending
suffix. Because $\beta$ is below the native spectrum, the combined order is
already ascending. At each depth compute the next prefix function once and try
the native suffix on it. If the current prefix initial value is negative,
stop extending that particular prefix: every longer ladder contains the same
failed test.

\par\noindent\begin{minipage}{\linewidth}
\begin{lstlisting}
enrichedLadder(f, maxExtra):
    if nativeLadder(f) succeeds: return it
    native := ascendingNativeFactors(f)
    for beta in selectedRatesBelow(min(native)):
        q := f
        prefix := []
        repeat at most maxExtra times:
            if q(0) < 0: break
            prefix.append(beta)
            q := (D - beta) q
            if nativeLadder(q) succeeds:
                return checkLadder(f, 0, prefix ++ native)
    return Unknown
\end{lstlisting}
\end{minipage}\par

The experiment used at most 96 auxiliary factors per selected rate. Its
results for \eqref{eq:epsilon-family} are:

\begin{center}
\begin{tabular}{@{}rrrrr@{}}
\toprule
$\varepsilon$ & Native result & $\beta$ & Extra factors & Total factors\\
\midrule
$1$    & unknown & $-4$  & 4  & 7\\
$1/2$  & unknown & $-8$  & 7  & 10\\
$1/4$  & unknown & $-16$ & 13 & 16\\
$1/16$ & unknown & $-64$ & 46 & 49\\
\bottomrule
\end{tabular}
\end{center}

These rows demonstrate a strict increase in the expressive power of the ladder
language. They are \emph{not} a superiority result over Forge's whole arithmetic
portfolio: an algebraic worker can establish the displayed square-plus-constant
identity immediately after introducing $u=e^x$. This family deliberately has a
known proof, making it a clean control for the enrichment mechanism.

The present search is not claimed complete for strictly positive functions on
a ray. The monotonicity proposition does not by itself show that the union of
all auxiliary-factor cones exhausts the strictly positive cone. Establishing
such a theorem, under explicit spectral and endpoint hypotheses, would be a
separate mathematical result. It must not be inferred from four successful
rows.

\subsection{An obstruction that no larger auxiliary budget repairs}

Set $\varepsilon=0$. The function $(e^x-2)^2$ is nonnegative on $[0,\infty)$,
is not identically zero, and vanishes at $x=\log2>0$. By
Theorem~\ref{thm:strict}, no finite ladder of the accepted form can establish
its nonnegativity on that entire ray, even with arbitrary real auxiliary
factors.

This is a structural failure, not evidence that the theorem is false. The
correct route is to preserve or discover the square factor and use algebraic
positivity. More generally, if the target factors as $h(x)^2r(x)$ and the
algebraic identity is proved, send $r$ to the analytic worker and keep the
square in the algebraic layer. The zeros of $h$ need not be located to prove
nonnegativity. To prove strict positivity, however, the source goal must also
supply or establish $h\ne0$ on the requested domain.

Another route is domain splitting at a known zero, but it needs a certified
split point and new initial-value arguments on each piece. A numerical root
approximation cannot be used as an exact domain boundary. Factoring is therefore
often substantially cheaper than asking a transcendental root isolator to
solve a problem that a square identity already settles.

\subsection{Towards better enrichment, without changing the checker}

The simple search exposes several concrete optimisations. Cache the prefix
polynomials so that increasing the number of copies does not recompute them.
Evaluate native-suffix initial values directly from jets before constructing
all suffix polynomials. Reject a proposed factor family using its earliest
negative coordinate. Use the coefficient norm and the minimum observed
positive coordinate to rank alternative $\beta$ values, but do not make that
heuristic part of acceptance.

A more general search can insert several distinct auxiliary rates. Theorem
\ref{thm:order} still removes their ordering dimension. What remains is a search
over multisets and their multiplicities. An LP relaxation could propose
candidate coordinates, but the nonlinear dependence on unknown rates must not
be described as a linear system. A conservative first implementation should
keep the exact discrete rate grid used here and extend it only when measured
misses justify the complexity.
